\section{Implementation}
\label{sec:impl}

\begin{guidance}[Implementation (15 marks)]
Choose \textbf{3 to 5 short code excerpts} that were \textbf{not straightforward}
to write, and explain each: what problem it solves, why you wrote it this way,
and what would go wrong with the obvious alternative. Good candidates:
receiving bytes in an interrupt and handing them to the main loop; the frame
assembler with timeout and oversize handling; non-blocking button debounce;
keeping the GUI responsive; separating unsolicited messages from replies.
Keep each listing short (10--25 lines) and narrow enough for one column; trim
comments and boilerplate, and say that you did. Do not paste whole files.
\end{guidance}

The listings in this report are trimmed: I removed comments, blank lines and
unrelated declarations.

\subsection{A register-level interrupt keeps up with 115200 baud}
Listing~\ref{lst:isr} is the receive path. At 115200 baud a byte arrives every
87\,$\mu$s. My first version re-armed \texttt{HAL\_UART\_Receive\_IT} after every
byte~\cite{um1785}. On the board it answered 6-character commands correctly,
but frames of 20 and 103 characters were answered with \texttt{NACK:BADCMD}
instead of \texttt{NACK:RANGE} and \texttt{NACK:TOOLONG}, which means the
server had flagged a UART fault mid-frame. I did not capture the overrun flag
directly, so overrun is my diagnosis from this evidence. The handler now reads
the status register and the data register itself~\cite{rm0091}, which is a few
dozen instructions long (I did not measure its run time). It clears the error flags, sets \texttt{rx\_fault}
so that the framer rejects the damaged frame, and pushes the byte into a
128-byte ring buffer. The interrupt only writes \texttt{rx\_head} and the main
loop only writes \texttt{rx\_tail}, so no lock is needed. The first version
also never enabled the USART1 interrupt in the NVIC, so no command arrived at
all; \texttt{HAL\_UART\_MspInit} now calls \texttt{HAL\_NVIC\_EnableIRQ}.

\begin{lstlisting}[float=tp, language=C, caption={Interrupt handler and ring buffer (\texttt{stm32f0xx\_it.c}, \texttt{main.c}).}, label=lst:isr]
void USART1_IRQHandler(void) {
    uint32_t isr = USART1->ISR;
    if (isr & (USART_ISR_ORE | USART_ISR_FE |
               USART_ISR_NE)) {
        USART1->ICR = USART_ICR_ORECF | USART_ICR_FECF
                    | USART_ICR_NCF;
        rx_fault = 1;               /* frame is damaged */
    }
    if (isr & USART_ISR_RXNE)
        uart_rx_push((uint8_t)USART1->RDR);
}
void uart_rx_push(uint8_t b) {
    uint16_t next = (rx_head + 1) % RX_RING_SIZE;
    if (next == rx_tail) rx_fault = 1;  /* ring full */
    else { rx_ring[rx_head] = b; rx_head = next; }
}
\end{lstlisting}

\subsection{The framer discards a bad line up to its terminator}
Listing~\ref{lst:framer} turns the byte stream into frames. The obvious
approach resets the buffer when a line is too long and carries on. Then the end
of a 100-character line would be parsed as a new command, so a line ending in
\texttt{NEW} would start a game. Instead the framer sets \texttt{frame\_err}
and drops every byte until the next terminator, then sends exactly one NACK.
A non-printable byte sets the same flag. A separate check in the main loop
compares \texttt{HAL\_GetTick()} with the time of the last received byte and
discards a half-received frame after 500~ms, so a lost terminator cannot glue
two commands together.

\begin{lstlisting}[float=tp, language=C, caption={Frame assembler (\texttt{main.c}).}, label=lst:framer]
while (rx_tail != rx_head) {
    char c = (char)rx_ring[rx_tail];
    rx_tail = (rx_tail + 1) % RX_RING_SIZE;
    last_rx_tick = HAL_GetTick();
    if (c == '\n' || c == '\r') {
        if (frame_err)
            send_str(frame_err == FE_TOOLONG
                ? "NACK:TOOLONG\n" : "NACK:BADCMD\n");
        else if (line_len > 0) {
            line_buf[line_len] = '\0';
            process_command(line_buf);
        }
        line_len = 0; frame_err = FE_NONE;
    } else if (frame_err) { /* drop until terminator */ }
    else if ((unsigned char)c < 0x20 ||
             (unsigned char)c > 0x7E)
        frame_err = FE_GARBAGE;
    else if (line_len < LINE_MAX_LEN)
        line_buf[line_len++] = c;
    else frame_err = FE_TOOLONG;
}
\end{lstlisting}

\subsection{The GUI waits on the port in a thread and draws from a queue}
Listing~\ref{lst:gui} shows how the client stays responsive (I split one line for width). A blocking
\texttt{readline} in a Tk callback would freeze the window for up to the
port timeout on every call. The worker thread is the only code that blocks on
the port; it hands lines to the Tk thread through a thread-safe queue, and the
Tk thread polls that queue every 50~ms and never blocks. The handshake uses
the same timer mechanism, so the GUI keeps sending \texttt{CONNECT} while the
board is still booting. Tkinter widgets are touched only from the Tk
thread~\cite{tkinter}.

\begin{lstlisting}[float=tp, language=Python, caption={Reader thread, queue polling and handshake (\texttt{gui\_client.py}).}, label=lst:gui]
def _read_thread(self, ser):       # worker thread
    while ser.is_open:
        raw = ser.readline()
        line = raw.decode(errors="replace").strip()
        if line:
            self.rx_queue.put(line)

def _poll_queue(self):             # Tk thread
    while not self.rx_queue.empty():
        self._handle_line(self.rx_queue.get())
    self.root.after(50, self._poll_queue)

def _handshake(self):              # until OK
    if not self.ser or self.connected: return
    if self.handshake_tries >= 15:
        self._close_port(); return
    self.handshake_tries += 1
    self._send("CONNECT")
    self.root.after(600, self._handshake)
\end{lstlisting}

\subsection{The AI searches a copy of the board and moves through the human's write path}
Listing~\ref{lst:ai} is the opponent's search: negamax with alpha-beta pruning,
scoring a win as $10-\mathit{depth}$ so that the AI prefers the quickest win and
the slowest loss. Two choices matter. First, the search places and removes marks
on a local copy of the board, so the game board changes only in
\texttt{apply\_move}, the single function that writes a mark; the human's
\texttt{MOVE} and the AI's reply both call it after their own checks, and the AI
re-checks that its chosen cell is empty first. A search that wrote to the real board
could leave a half-explored position behind. Second, cells are tried centre,
corners, then edges, which cuts the search: on a PC build of the same source the
worst single decision, the first reply, takes 1\,260 calls to \texttt{negamax}
(mean 35), against up to $8! = 40\,320$ move sequences without pruning. I have
not timed the search on the board, but in all five full test runs every AI reply
arrived inside the harness's 2~s read timeout. The search runs in the main loop
while the receive interrupt keeps buffering bytes.

\begin{lstlisting}[float=tp, language=C, caption={AI search (\texttt{main.c}, trimmed).}, label=lst:ai]
int negamax(char *b, char me, char opp,
            int depth, int alpha, int beta) {
    int best = -100;
    for (int k = 0; k < 9; k++) {
        int i = ai_order[k], sc;   /* 4,0,2,6,8,1,3,5,7 */
        if (b[i] != '.') continue;
        b[i] = me;
        if (line_win(b, me))    sc = 10 - depth;
        else if (board_full(b)) sc = 0;
        else sc = -negamax(b, opp, me, depth + 1,
                           -beta, -alpha);
        b[i] = '.';
        if (sc > best)     best = sc;
        if (best > alpha)  alpha = best;
        if (alpha >= beta) break;
    }
    return best;
}
\end{lstlisting}
